\section{Population, Contracts and Experimental Design}
\label{sec:methodology}
\subsection{Population and evidence units}
Table~\ref{tab:final-population} states the final denominator. Historical 300-pass labels comprised 259 native-check and 41 stdout comparisons; they were selection evidence, not proof of identical-input array equivalence. Source audits and new contracts replace that inference. A benchmark name can contain several configurations and qualified modes; one name does not certify every CLI, precision or horizon.
\begin{table}[t]\centering\small
\caption{Completed assessment: 447 distinct names. Numerical contracts include corrected-reference, known-defect and revised-acceptance outcomes; agreement is not universal scientific validity.}\label{tab:final-population}
\begin{tabular}{lr}\toprule Outcome & Names\\\midrule
Numerical contracts &436\\Property-only contracts &2\\Execution-only contracts &5\\Native hardware barriers &4\\\midrule Total assessed &447\\Pending &0\\\bottomrule\end{tabular}
\end{table}

Inputs, shapes, storage/computation precision and tolerances are declared per configuration. Per-record input evidence includes uploaded fixtures, exact prefix comparisons and source/buffer reviews; independent initial-device readback is not universal. Every accepted trial checks its recorded outputs against both GPU references; independent analytic, higher-precision, invariant or discrete witnesses provide additional guarantees where recorded. Full state, aggregate outputs, canonicalized results and property scores remain different evidence strengths. Guards, nonfinite handling, tolerance units and inspected counts are explicit. A local PASS cannot override a failed reference comparison.

Reference status is separate from agreement kind. The conservative per-model census has 383 ordinary named contracts (including two property cases), 42 corrected-native, 12 known-native-defect, one revised-acceptance, five execution-only and four hardware-blocked names. Where a name has both untouched and corrected modes, performance eligibility is decided per configuration. The four unavailable vendor operations are FP4 GEMM, FP8 GEMM, MXFP6 GEMM and MXFP8 GEMM on the recorded environment; they are not demonstrated KA/JACC language limitations.

\subsection{Scientific intent and implementation descriptors}
Scientific intent records the operator, domain, numerical method, observable, accuracy and scale. Implementation records describe actual algorithms, mapping, layout, arithmetic, synchronization, library use and backend dependence. These are operational JSON relations, not automated logical inference. Source-backed labels are distinct from profiled properties. The revision introduces a controlled implementation vocabulary while preserving historical narrative labels and ambiguous/unknown fields. Configuration overrides take precedence over benchmark-wide descriptions. Explicit timing-boundary categories cover 81 of 428 eligible configurations; the remaining raw scope descriptions are retained unreconciled. Extractable numeric dimensions cover 173 configurations but have incompatible units and are not pooled as one size axis.
\input{sections/generated/revision_classes}

Scientific families support discovery; they do not identify bottlenecks. For example, soil-law updates require admissible constitutive inputs, wave operators require retained time levels and boundaries, and stochastic components require an explicit stream or distribution contract. The artifact's executable query links those requirements to implementation classes and measured configurations rather than inferring capability from a directory name.

\subsection{Timing and environment}
Every primary model/configuration has five checked trials; the prescribed mean is
\[\bar t=(t_2+t_3+t_4+t_5)/4,\qquad S_a=\bar t_{\rm CUDA}/\bar t_a.\]
State is reset, setup/JIT and output capture are outside the declared region, and matched completion boundaries are checked. Some regions include allocation, transfers, CPU control or library calls; they are not pooled as isolated device latency. Unchanged reference series can be reused only with dependency/protocol checks. All samples, including spikes, remain retained.

Single-device measurements use a Tesla V100 PCIe 32~GB, compute capability 7.0, driver 580.178.04, Julia 1.12.6, CUDA.jl 6.2.1, KA 0.9.45 and JACC 1.4.0. CUDA 12.0 is the primary native toolkit; recorded exceptions use 12.9. Runtime library identities are separate from link arguments. Special-function environments add SpecialFunctions 2.9.0. Native builds use \texttt{-O3} and \texttt{sm\_70} with per-case arithmetic flags. Two-V100 configurations are separate from single-device aggregates. Nsight Compute 2022.4.1 and a separately calibrated 2025.2.0 recovery stratum profile selected kernels; method identities remain explicit. Profiler rates are not pooled across methods, and primary timing series are unchanged.

The original shared-host measurements do not fix clocks. New independent five-trial blocks counterbalance model order and retain clock/load telemetry; they supplement, rather than replace, the frozen results. For family summaries, each name contributes the geometric mean of its configuration ratios; family markers are medians with interquartile spread across names. Those spreads are not timing confidence intervals. Profiled replay duration and host-inclusive execution are different metrics.

\paragraph{Prospective V100 counter campaign.}
A separate, currently incomplete campaign assesses all 537 native names and four
versions, retaining native CUDA as the baseline and the frozen 447-name results
and 90-name extension as distinct histories. Each profile binds the selected
configuration, executed source/library, device and kernel geometry. Nsight Compute
2022.4.1 uses cache control \texttt{all} and clock control \texttt{none}; ceilings
are architectural work per cycle scaled by observed clocks, not measured sustained
limits. Unlike Figure~\ref{fig:triad-roofline}'s useful-work numerator, new points
use executed add/multiply/FMA counters separately by precision and measured DRAM
read-plus-write bytes. Packed-half and tensor FLOP conversions require validation;
otherwise those components remain unknown.
